\section{Workflow with pretrained transformer models}
\label{sec:pretrained}

The three workflows in Figure \ref{fig:workflow} handle the text data differently.

In the traditional workflow, the result of text preprocessing would enable one to calculate
a similarity measure of any given record pair $(i,j)$ with associated texts
$(\tau_i,\tau_j)$. The Jaro--Winkler metric
\citep{jaro1989advances,winkler1990string} is very common in practice, which yields a
value between 0 and 1 if appropriately configured. It captures the lexical overlaps, where
each text is treated as a string of characters, but is not effective for expressions that differ
only superficially to the human eyes, such as ``Conf. Proc.'' and ``Proceedings''.

In the proposed workflow, preprocessing is replaced by embedding, which uses a bi-encoder
model to turn each given text $\tau_i$ into a numerical vector, denoted by $\mathbf{x}_i$,
which typically has a dimension of hundreds or more. A common similarity measure of
$(\tau_i,\tau_j)$ is the cosine similarity of $(\mathbf{x}_i,\mathbf{x}_j)$, given as

\begin{equation}
\gamma_{ij}
=
\frac{\mathbf{x}_i^\top \mathbf{x}_j}
{\lVert \mathbf{x}_i\rVert \lVert \mathbf{x}_j\rVert}.
\end{equation}

The bi-encoder similarity measure can capture semantic similarity that is not
represented by character-level overlap, while the embeddings can be calculated
efficiently and used routinely for the semantic blocking needed before
scoring-classification.

In the workflow using a cross-encoder model, which takes a pair of texts $\tau_i$ and
$\tau_j$ as input directly, one can obtain a score for the similarity between them. A cross-encoder can represent pair-specific interactions between the two texts
that similarity measures calculated from separate embeddings, such as
$\gamma_{ij}$ above, cannot represent directly. Using a
bi-encoder for semantic blocking, and applying a cross-encoder to the candidate-paired
texts can potentially achieve the best link classification performance.

Table \ref{tab:silverman-example} illustrates the differences with two citations of
Silverman's 1986 monograph. The Jaro--Winkler metric is only calculated for comparisons
per field. The pretrained bi-encoder model \texttt{all-MiniLM-L6-v2}
\citep{reimers2019sentence} yields 384-dimensional embeddings, which result in the cosine
similarity measures here. The pretrained cross-encoder model
\texttt{ms-marco-MiniLM-L-6-v2} \citep{wang2020minilm} yields scores from paired texts
directly without numerical embeddings being compared afterwards.

\begin{table}[ht]
\centering
\caption{Per-field and whole-record scores for two citations of Silverman's 1986
monograph: ``Silverman, B. W. (1986). Density Estimation for Statistics and Data
Analysis. Chapman and Hall, London.'' and ``Bernard W. Silverman (1986). Density
estimation for statistics and data analysis. Monographs on Statistics and Applied
Probability.''}
\begin{tabular}{lccc}
\toprule
Text & Jaro--Winkler metric & Bi-encoder cosine & Cross-encoder score \\
\midrule
Author        & 0.510 & 0.784 & 0.974 \\
Title         & 0.856 & 0.872 & 0.991 \\
Venue         & 0.933 & 0.969 & 0.999 \\
Year          & 1.000 & 1.000 & 0.999 \\
Full citation & --    & 0.887 & 0.999 \\
\bottomrule
\end{tabular}
\label{tab:silverman-example}
\end{table}

All methods rank the field agreements in the same order, from the lowest on author
name to the highest on year. However, the numerical values produced by the three methods
should not be compared as if they were measured on a common scale: the Jaro--Winkler
metric, cosine similarity and cross-encoder output arise from different scoring rules. The
example instead illustrates the difference in representation. The cross-encoder evaluates
the two names jointly, whereas the bi-encoder must commit ``Silverman, B. W.'' to a fixed
vector without the sight of ``Bernard W. Silverman'', and vice versa. Joint encoding allows
pair-specific interactions to be represented directly, which separate embeddings cannot do.
Whether this architectural advantage translates into lower linkage error is assessed below
using the FLR and MMR.

Notice that, for deterministic record linkage, one can already apply a chosen threshold
to any of the similarity measures for link classification directly. The cross-encoder has greater flexibility for representing pair-specific
interactions, but the computation cost of applying it may be too high for many
applications of entity resolution even with bi-encoder blocking.
Below we first consider how the proposed workflow of probabilistic linkage allows one to
incorporate the traditional models for scoring-classification, so as to reduce the
computational cost, albeit without fine-tuning of the transformer models; the workflow
with fine-tuning of transformer models will be considered in Section
\ref{sec:finetuning} later.


\subsection{Embedding and scoring}
\label{sec:embedding-scoring}

For scoring by generative or discriminative models, one needs numerical representation of
key-value comparisons, denoted by $\gamma_{h,ij}$ for each field $h$ of $\tau_i$ and
$\tau_j$, where $h=1,\ldots,H$. The Jaro--Winkler metric is a possibility for the
traditional workflow, whereas one can use the cosine similarity measure of bi-encoder
embeddings for the proposed workflow.

By the MEC approach using generative models, one would obtain an estimated density
ratio for the record pair $(i,j)$, which is given as

\begin{equation}
\label{eq:density-ratio}
r_{ij}
=
\prod_{h=1}^{H}
\frac{
f(\gamma_{h,ij}\mid y_{ij}=1)
}{
f(\gamma_{h,ij}\mid y_{ij}=0)
}.
\end{equation}

\rev{The product form in Equation \eqref{eq:density-ratio} assumes that the
field-level comparison outcomes are conditionally independent given the match
status $y_{ij}$. This assumption may be violated for related bibliographic fields,
such as title, authors and venue. In that case, the product of the field-specific
density ratios need not equal the likelihood ratio based on their joint distribution.}

Notice that discrete comparison outcomes, i.e.\ agreement or not, are common for record
linkage traditionally, in which case one would replace the density ratio by a probability
ratio. Given continuous comparison outcomes $\gamma_{h,ij}$, we simply apply kernel
density estimation \citep{silverman1986density} to estimate the two densities in
Equation \eqref{eq:density-ratio}, separately, in the training set of matches and non-matches that are constructed as
described in Section \ref{sec:training-candidates}.

Meanwhile, discriminative models target the probability of match given the key-value
comparison outcomes, denoted by

\begin{equation}
\label{eq:match-probability}
p_{ij}
=
\Pr\left(
y_{ij}=1
\mid
\gamma_{1,ij},\ldots,\gamma_{H,ij}
\right).
\end{equation}

We use logistic regression as a basic discriminative model. Other models are possible,
such as support vector machine (SVM), gradient-boosted trees (XGBoost) and multilayer
perceptron (MLP) neural network. A given discriminative model can be estimated from
the training set of matches and non-matches as a single dataset.

Using the labelled records in the DBLP-Scholar and Abt-Buy datasets, we consider
blocking scheme-II (Section 1.3) with fixed size $k$ for each record in $A(s_M)$. This
induces a floor MMR, since a file-$A$ record cannot be linked correctly when none of its
valid matched records is included in the candidate block.

\begin{table}[ht]
\centering
\caption{Floor MMR by block size $k$ in DBLP-Scholar and Abt-Buy datasets.
Floor-zero $k$ refers to the minimum block size needed for floor MMR to be zero.}
\begin{tabular}{lrrrr}
\toprule
Dataset & $k=1$ & $k=5$ & $k=50$ & Floor-zero $k$ \\
\midrule
DBLP-Scholar & 0.022 & 0.006 & 0.001 & 4535 \\
Abt-Buy      & 0.278 & 0.082 & 0.003 & 748 \\
\bottomrule
\end{tabular}
\label{tab:floor-mmr}
\end{table}

Table \ref{tab:floor-mmr} shows the floor MMR for the two datasets, where the candidates
for each record in $A(s_M)$ are selected from $B$ based on the bi-encoder cosine
similarity. At $k=1$, it is 0.022 for the former and 0.278 for the latter, which means
that about 97.8\% of the matchable file-$A$ records have a valid match as the top
candidate in the DBLP-Scholar dataset, whereas about 72.2\% do so in the Abt-Buy
dataset. This indicates that entity resolution is more difficult in the latter case.

Meanwhile, both the floor MMRs become practically acceptable once $k$ is increased to
50, especially because one would have to set $k=4535$ to remove the floor (i.e.\ making
the floor MMR zero) for the DBLP-Scholar dataset, or $k=748$ for the Abt-Buy dataset.


\subsection{Comparison of embedding and scoring methods}
\label{sec:comparison-pretrained}

The comparison outcomes $\gamma_{h,ij}$ are calculated for the candidate pairs
constructed as described in Section \ref{sec:training-candidates}, either as the
Jaro--Winkler metric or the cosine similarity of bi-encoder embeddings. Given
$\{\gamma_{h,ij}\}$, we compare the resulting link classification errors under
generative and discriminative scoring models. The FLR and MMR are calculated on
the held-out file-$A$ records according to Section \ref{sec:evaluation}.

The distinction between the two scoring approaches is worth explaining. For MEC
under the generative model, we link each file-$A$ record $i$ to the candidate
$j\in C_i$ with the highest density ratio $r_{ij}$ in Equation
\eqref{eq:density-ratio}. Thus, the best available candidate is always selected.
For a discriminative model one could similarly select the candidate with the
highest match probability $p_{ij}$ in Equation \eqref{eq:match-probability}.
However, this would force a link even when all the candidate probabilities are
small.

We therefore consider a threshold-based classification rule for the discriminative
models. As described in Section \ref{sec:training-candidates}, a validation subset
of the training pairs is used to choose the threshold that maximises the F1 score.
A link is declared only when the highest-scoring candidate exceeds this threshold;
otherwise the file-$A$ record is left unlinked. This provides a useful contrast
with the generative procedure because it allows the model to abstain when the
candidate evidence is weak. We do not investigate alternative cross-validation
schemes or decision criteria here, since our purpose is to illustrate the role of
discriminative scoring-classification within the proposed workflow rather than to
optimise the procedure for these two datasets.

\begin{table}[ht]
\centering
\caption{FLR and MMR by key-value comparison method, scoring-classification model
and block size $k$. DBLP-Scholar (top half), Abt-Buy (bottom half).
\rev{Entries give the mean (Monte Carlo standard error) over three random splits.}}
\scriptsize
\setlength{\tabcolsep}{3pt}
\begin{tabular}{llrrrr}
\toprule
& & \multicolumn{2}{c}{$k=5$} & \multicolumn{2}{c}{$k=50$} \\
\cmidrule(lr){3-4}\cmidrule(lr){5-6}
Key-value comparison & Scoring-classification & FLR & MMR & FLR & MMR \\
\midrule
String, Jaro--Winkler & Generative, kernel density & 0.034 \rev{(0.0026)} & 0.018 \rev{(0.0005)} & 0.046 \rev{(0.0031)} & 0.013 \rev{(0.0002)} \\
 & Discriminative, logistic & 0.030 \rev{(0.0038)} & 0.036 \rev{(0.0030)} & 0.027 \rev{(0.0027)} & 0.050 \rev{(0.0034)} \\
Bi-encoder, cosine & Generative, kernel density & 0.023 \rev{(0.0013)} & 0.018 \rev{(0.0015)} & 0.043 \rev{(0.0026)} & 0.014 \rev{(0.0009)} \\
 & Discriminative, logistic & 0.016 \rev{(0.0038)} & 0.056 \rev{(0.0028)} & 0.009 \rev{(0.0027)} & 0.068 \rev{(0.0027)} \\
\midrule
String, Jaro--Winkler & Generative, kernel density & 0.389 \rev{(0.0089)} & 0.500 \rev{(0.0043)} & 0.449 \rev{(0.0097)} & 0.525 \rev{(0.0075)} \\
 & Discriminative, logistic & 0.343 \rev{(0.0279)} & 0.597 \rev{(0.0385)} & 0.338 \rev{(0.0081)} & 0.662 \rev{(0.0468)} \\
Bi-encoder, cosine & Generative, kernel density & 0.278 \rev{(0.0044)} & 0.425 \rev{(0.0062)} & 0.327 \rev{(0.0043)} & 0.356 \rev{(0.0073)} \\
 & Discriminative, logistic & 0.248 \rev{(0.0191)} & 0.479 \rev{(0.0505)} & 0.233 \rev{(0.0210)} & 0.543 \rev{(0.0597)} \\
\bottomrule
\end{tabular}
\label{tab:comparison}
\end{table}

Table \ref{tab:comparison} reports the test set FLR and MMR for four combinations of
embedding and scoring methods, given two block sizes $k$. The results are averaged over
\rev{three random 50--50 training-test splits, with Monte Carlo standard errors in
parentheses. Comparisons of close means should be read in light of this
split-to-split variation.}

Regarding bi-encoder embeddings instead of the traditional character-string comparisons,
the effects are not obvious for the DBLP-Scholar dataset where the texts are less messy
(Table 1). The bi-encoder yields lower FLRs in all four comparisons, but does not
consistently reduce the MMR. In contrast, the improvements are clearly visible for the
Abt-Buy dataset with more messy texts. Both the FLR and MMR are reduced by the
bi-encoder representation for either scoring model and either block size.

These results illustrate the potential advantages of transformer-based
embeddings for text-based entity resolution. Notice that the performance of the bi-encoder can be further changed
by fine-tuning, which will be considered in Section \ref{sec:finetuning}.

Regarding scoring-classification by generative or discriminative models, the results here
show that the discriminative logistic regression model yields lower FLRs than the
generative kernel density model, whereas the latter yields lower MMRs. The distinction
is especially clear for the Abt-Buy dataset. This reflects the different link-classification
rules used here: we always link the top-ranked candidate under the generative model,
whereas a validation split is used to determine a threshold for the discriminative model,
which may abstain from declaring a link.
\rev{Accordingly, the FLR/MMR contrast in
Table \ref{tab:comparison} should not be interpreted as an intrinsic difference between
generative and discriminative model families; it also reflects the different decision rules
applied to their scores.}

Without advocating for one modelling approach against the other generally, we note that
scoring-classification by discriminative models is worth considering in practice, although
such models have been less common in record linkage traditionally.

For completeness, we also apply the pretrained cross-encoder directly to the candidate
pairs obtained from bi-encoder blocking. Since the cross-encoder operates on paired
record texts rather than on the numerical comparison outcomes used by the traditional
scoring models, we report its results together with those of the fine-tuned cross-encoder
in Table \ref{tab:fine-tuning}. This allows the effect of task-specific fine-tuning to be
assessed separately from that of using the cross-encoder architecture itself.

Regarding the proposed workflow, using the bi-encoder model for embedding and
the traditional generative or discriminative models for automated
scoring-classification can reduce the number of candidate pairs that need to be
processed by a cross-encoder, compared with applying the cross-encoder to every
candidate block. This provides the basis for examining the accuracy-cost tradeoff
after fine-tuning, compared with either the traditional workflow without transformer
models or the transformer-only workflow.

\subsection{Sensitivity analysis}
\label{sec:sensitivity}

Here we conduct sensitivity analysis of a methodological choice not covered above, which
exemplifies the relevant considerations in practice and helps to prepare the ground for
fine-tuning transformer models later.

\begin{table}[ht]
\centering
\caption{FLR and MMR by different discriminative models. Entries give the
mean (Monte Carlo standard error) over three random splits.}
\scriptsize
\setlength{\tabcolsep}{4pt}
\begin{tabular}{llrrrr}
\toprule
& & \multicolumn{2}{c}{$k=5$} & \multicolumn{2}{c}{$k=50$} \\
\cmidrule(lr){3-4}\cmidrule(lr){5-6}
Dataset & Model & FLR & MMR & FLR & MMR \\
\midrule
DBLP-Scholar & Logistic & 0.016 \rev{(0.0038)} & 0.056 \rev{(0.0028)} & 0.009 \rev{(0.0027)} & 0.068 \rev{(0.0027)} \\
 & SVM & 0.017 \rev{(0.0021)} & 0.040 \rev{(0.0037)} & 0.018 \rev{(0.0017)} & 0.041 \rev{(0.0048)} \\
 & XGB & 0.033 \rev{(0.0030)} & 0.021 \rev{(0.0014)} & 0.024 \rev{(0.0055)} & 0.041 \rev{(0.0050)} \\
 & MLP & 0.030 \rev{(0.0025)} & 0.030 \rev{(0.0016)} & 0.026 \rev{(0.0024)} & 0.034 \rev{(0.0066)} \\
\addlinespace
Abt-Buy & Logistic & 0.248 \rev{(0.0191)} & 0.479 \rev{(0.0505)} & 0.233 \rev{(0.0210)} & 0.543 \rev{(0.0597)} \\
 & SVM & 0.252 \rev{(0.0170)} & 0.592 \rev{(0.0274)} & 0.211 \rev{(0.0235)} & 0.701 \rev{(0.0348)} \\
 & XGB & 0.399 \rev{(0.0096)} & 0.538 \rev{(0.0187)} & 0.296 \rev{(0.0130)} & 0.551 \rev{(0.0235)} \\
 & MLP & 0.251 \rev{(0.0111)} & 0.438 \rev{(0.0238)} & 0.228 \rev{(0.0202)} & 0.565 \rev{(0.0704)} \\
\bottomrule
\end{tabular}
\label{tab:discriminative-sensitivity}
\end{table}

Instead of logistic regression, one can use many other discriminative models. Table
\ref{tab:discriminative-sensitivity} compares it to a radial-kernel SVM, an XGBoost and
a small MLP net, all based on the same bi-encoder embeddings. The FLR and MMR are
obtained in the same way as above, and the results of logistic regression are reproduced
from Table \ref{tab:comparison} for easy comparison. The models exhibit different
tradeoffs between the two error measures. \rev{For DBLP-Scholar, logistic regression has
low FLR in both block-size configurations, while several models have lower observed
MMR. The FLR estimates for logistic regression and SVM at $k=5$ are close relative
to their Monte Carlo standard errors. For Abt-Buy, the FLR and MMR comparisons vary
with the model and block size, with appreciable split-to-split variation in several
entries. Logistic regression provides a simple discriminative model and will be used
as the default for the investigation later.}